\section{Evaluation}
\label{sec:eval}

The evaluation answers four research questions. \textbf{RQ1}: do the formal
constraints guarantee experiment integrity, and does the verification detect
injected violations? \textbf{RQ2}: does the plugin architecture support
independent extension without kernel changes? \textbf{RQ3}: are experiments
traceable and reproducible? \textbf{RQ4}: does the framework sustain
heterogeneous data and acceptable overhead? Unless stated otherwise, all
numbers come from artefacts committed with the repository (TLC log,
detection-matrix JSON, result bundles, benchmark scripts).

\subsection{RQ1: Integrity verification}

\paragraph{Model checking.}
TLC explores the lifecycle model exhaustively over the configured finite
constants and reports no errors: 89{,}476 states generated, 46{,}686 distinct
states, depth 11, no deadlocks, all nine properties (type safety, INV-1
through INV-6, fail-closed discipline) holding in every reachable state.
Every rejection path (tampered partitions, foreign-partition windows,
test-scope fits, contract-less plugins, incomplete bundles) is exercised by
construction, because the spec offers the bad candidate non-deterministically
and TLC explores it.

\paragraph{Property-based testing.}
Eight oracle properties run in the default suite. Each pairs a generated
input with an \emph{independently recomputed} expected outcome---for
example, the recording audit must pass if and only if the generated set has
unique identifiers, supported labels, and fully assigned conditions---so
agreement with the oracle is asserted, not merely acceptance of valid cases.
The properties cover task-partition validation, recording-audit agreement,
window-partition totality and exclusivity under relabelling, prepared-table
round trips, window-identity uniqueness, seed bijection, training-policy
validation, and plugin-contract rejection.

\paragraph{Fault injection.}
Table~\ref{tab:mutation} shows the twelve-entry catalogue: every planted
violation is rejected fail-closed by at least one mechanism, with the
violated invariant named in the error. Nine distinct mechanisms participate.
The matrix reporter itself is tested for honesty: a canary test disables one
guard and requires the report to say \emph{undetected} rather than silently
passing.

\begin{table}[t]
\centering\small
\caption{Fault-injection detection matrix (all entries executed against the
live kernel).}
\label{tab:mutation}
\begin{tabular}{@{}llp{0.46\linewidth}@{}}
\toprule
Mutation & Target & Detecting mechanism \\
\midrule
M-LEAK-01 & INV-2 & pre-training recording audit \\
M-LEAK-02 & INV-2 & recording cross-audit vs.\ canonical metadata \\
M-LEAK-03 & INV-3 & kernel fit-scope guard (state digest differs from all-rows fit) \\
M-LEAK-04 & INV-3 & partition-restricted views (fit never sees test rows) \\
M-HASH-01 & INV-1 & adapter checksum verification at open \\
M-SEED-01 & INV-4 & result validator (bijection) \\
M-SEED-02 & INV-4 & result validator (duplicate seed) \\
M-PROV-01 & INV-5 & result validator (required fields) \\
M-PROV-02 & INV-5 & input-artifact hash recheck after tampering \\
M-PLUGIN-01 & INV-6 & plugin registry identity check \\
M-PLUGIN-02 & INV-6 & plugin contract validation \\
M-PLUGIN-03 & INV-6 & plugin registry collision rejection \\
\bottomrule
\end{tabular}
\end{table}

The enforcement is not decorative: during the gearbox integration, the
length-consistency check at stream open rejected the first real extraction
run because the registry's row counts---naively counted---overstated each
file by two trailing blank lines. The registry was corrected before any
result existed, and the incident is preserved in the integration decision
record.

\subsection{RQ2: Independent extension}

Table~\ref{tab:integration} quantifies the two most recent integrations. All
adapter-side code, tests, registries, task manifests, and documentation sum
to 812 and 595 lines respectively; the lifecycle kernel gained no
dataset-specific logic. A committed test enforces this structurally: no
kernel module may contain a dataset identifier, with the two known debts
(the fetch dispatcher and one frozen demonstration asset) declared in the
test itself. Kernel touches that did occur were generic (an optional
sampling-rate field for sources that state none) and are classified in the
integration-effort audit.

\begin{table}[t]
\centering\small
\caption{Integration effort for the two most recent datasets.}
\label{tab:integration}
\begin{tabular}{@{}lrrr@{}}
\toprule
Integration & Adapter LOC & Test LOC & Files added (all artefacts) \\
\midrule
Gearbox (20 recordings) & 239 & 175 & 6 \\
KAIST acoustic (5 recordings) & 261 & 158 & 5 \\
\bottomrule
\end{tabular}
\end{table}

Compared with the script-per-dataset tradition, the operations that used to
require editing experiment code---adding a dataset, adding a model, changing
the held-out conditions, checking leakage, reproducing a result---now live
in, respectively, an adapter package, a model plugin, a task manifest, an
automated audit, and a result bundle with hash recheck. The boundary is
candid: both audited integrations were performed by the author; soliciting
one non-author integration is an open item, and the repository ships an
issue template with an acceptance checklist for exactly that purpose.

\subsection{RQ3: Traceability and reproduction}

Traceability is content-addressed: every result bundle hashes its task
manifest, feature source (table or dataset manifest), and training policy,
and \texttt{results validate -{}-check-hashes} recomputes them; tampering
with an input after the fact is detected as a mismatch (M-PROV-02). The
strongest reproduction evidence to date is same-commit, cross-Python: the
development environment (3.11) and an independent virtual environment (3.12,
package and external plugin installed from source) produce identical
per-seed test Macro-F1 and worst-condition scores on a frozen task. A CI
gate extends this across operating systems---the same task runs on Ubuntu
and macOS runners and the comparison job fails on any divergence in
canonical metrics; its first executed run accompanies the public push of
this repository. Full multi-machine, multi-investigator reproduction
remains future work.

\subsection{RQ4: Heterogeneity, scale, and overhead}

\paragraph{Heterogeneous tasks, executed on real data.}
Three task families newly enabled by the v1.2 architecture were executed on
real archives end to end (streaming from checksum-verified sources, three
frozen seeds each, group-bootstrap intervals over test recordings). Under a
deliberately reduced two-epoch \emph{audit policy}---chosen to validate
mechanism, not to claim scores---the gearbox cross-load task yields
window-level Macro-F1 of $0.333$ with a recording-level 95\% CI of
$[0.11, 0.47]$ over eight test recordings; the combined-condition
load$\times$defect-size task yields $0.330$--$0.337$ with CIs around
$[0.14, 0.46]$ over six test recordings. The five-fold leave-one-device-out
family on the bearing archive produces fold means of
$0.565/0.377/0.597/0.472/0.268$ (outer mean $0.456$) with large per-seed
variance---consistent with, and sharper than, the frozen v1.0 finding that
device hold-out fails badly while load hold-out looks deceptively good.
These runs demonstrate what the framework now makes routine---combined
shifts, device rotation, and recording-level uncertainty on real,
hash-verified archives---and are labelled non-protocol in the artefacts; the
full-protocol v1.0/v1.1 matrices (3 tasks $\times$ 4 models and 2 tasks
$\times$ 4 models, with negative results retained) were produced by earlier
versions of the same pipeline and remain the benchmark claims of record.

\paragraph{Framework overhead.}
Table~\ref{tab:overhead} measures the kernel's cost on synthetic tables.
Loading plus schema-validating a 100{,}000-row v2 table costs 8.8\,s against
0.57\,s for bare JSON parsing ($15.6\times$)---the price of constructing
immutable, identity-carrying rows and validating them; the optional Parquet
backend addresses exactly this cost for large matrices. Result validation is
0.7\,ms per bundle. Streaming extraction plus audit of the 36-file triaxial
archive (1{,}154 windows) completes in 1.4\,s.

\begin{table}[t]
\centering\small
\caption{Kernel load-and-validate cost versus bare \texttt{json.loads}
(median of five repeats).}
\label{tab:overhead}
\begin{tabular}{@{}rrrr@{}}
\toprule
Rows & Kernel (s) & Bare JSON (s) & Ratio \\
\midrule
1{,}000 & 0.086 & 0.0028 & 30.1 \\
10{,}000 & 0.780 & 0.0287 & 27.1 \\
100{,}000 & 8.813 & 0.5667 & 15.6 \\
\bottomrule
\end{tabular}
\end{table}

\paragraph{Modality and format coverage.}
The signal layer presently streams four source formats (MAT v5, Siemens
Test.Lab MAT, CSV, tab-separated text) across vibration, acoustic pressure,
and gearbox accelerometers, with per-file rates from 25.6\,kHz to 51.2\,kHz
and one archive whose rate is honestly \texttt{None}. One archive (the
multimodal vibration set, 45 files) is integrated and audit-verified but its
full end-to-end run awaits a machine with more memory than the development
host (8\,GB); the data card records this boundary.

\subsection{Threats to validity}

\emph{Internal.} The fault catalogue plants artefact-level violations by
hand; it is not exhaustive source-code mutation, and a guard outside the
catalogue's reach is untested until a property or real incident exercises
it. \emph{External.} Verification uses many synthetic fixtures; real-archive
coverage exists for four of six integrations, and the audit-policy training
runs are not comparable to the frozen protocol results. \emph{Construct.}
``Formally specified'' covers the abstract lifecycle; implementation claims
rest on tests, runtime checks, and injection results, and the paper keeps
the two layers separate. \emph{Adoption.} All integrations and audits to
date are author-performed; third-party evidence is solicited but absent.
